\printedpage{14}
% head: A. Grothendieck : Dérivateurs, Ch. XIX, Édition des Universités de Montpellier II et Paris VII
% foot: Version 11 février 2019

\msp{page 23}
The decomposition axiom is satisfied for every prederivator of the form $\mathbf{D}_{\mathcal{M},W}$. The left exactness axiom Der 5 bis is valid, provided that the $(\mathcal{M}(I),W(I))$ fit into a structure with cofibrations, which is the case. Likewise for the localization axiom Der 2, provided that in addition $W$ is strongly saturated, which is also the case. Finally, the final object axiom says that if $I$ has a final object, then
\[ \Ho_W(e) \longrightarrow \Ho_W(I) \]
is fully faithful, or again that the adjunction morphism
\[ u_!u^*(\xi) \longrightarrow \xi \]
is an isomorphism. But by virtue of 1.12, case 3°, in this case $u^*$ maps $\mathcal{M}(e)_{\mathrm{c}} = \mathcal{M}_{\mathrm{c}}$ into $\mathcal{M}(I)_{\mathrm{c}}$, so the computation of the pair of adjoint functors $u_!^W$, $u_W^*$ can be carried out using the restrictions of $u_!$, $u^*$ to $\mathcal{M}(I)_{\mathrm{c}}$ and $\mathcal{M}_{\mathrm{c}}$ respectively. Since the adjunction morphism $u_!^{\mathrm{c}}u_{\mathrm{c}}^* \longrightarrow \mathrm{id}$ is an isomorphism, the same therefore holds for $u_!^W u_W^* \longrightarrow \mathrm{id}$, which is deduced from it by localizing, q.e.d.

\textbf{2.5.} To finish constructing $\Ho_W$ as a derivator on all of Cat, one must \msp{page 24} therefore establish Der 4, Der 5: existence of the $u_*^W$, and computation of the fibers of $u_*^W$. As for Der 6, it is obtained automatically, from the fact that each $(\mathcal{M}(I),W(I))$ fits into a fibration structure $(W(I),F(I))$. (Omission in the proof of 2.4.) The natural idea is now to use the triple structure
\[ (W(I),C(I),\ \underbrace{F[I]}_{\deq TC(I)_*}\ )\ , \tag*{(2.5.1)} \]
but the difficulty is to prove the factorization for a couple
\[ C(I) \longleftrightarrow TF[I]\ , \qquad \underbrace{TC(I)}_{\deq C(I)\cap W(I)} \quad \longleftrightarrow F[I]\ , \tag*{(2.5.2)} \]
for want of being able to show, under the general conditions of 2.3 alone (special Quillen triple) that $C(I)$ and $TC(I)$ are also special. (Cf. question 1.11.) But with suitable conditions on $\mathcal{M}$, we shall manage to get by if $C = \mathit{mono}$. The relevant statement, a little technical, is found in XV cor. 7.4 (p. 98-99), which I am going to recall here.

\printedpage{15}
% head: A. Grothendieck : Dérivateurs, Ch. XIX, Édition des Universités de Montpellier II et Paris VII
% foot: Version 11 février 2019

\msp{page 25}
\textbf{Theorem 2.6.} \textit{Let $\mathcal{M}$ be a category satisfying the following conditions.}
\[ \left\{ \begin{minipage}{0.72\textwidth}\itshape
\begin{itemize}[leftmargin=2.5em]
\item[\upshape a)] $\mathcal{M}$ is ``special'', i.e. stable under small $\varinjlim$ , and its elements are accessible.
\item[\upshape b)] Filtered inductive limits in $\mathcal{M}$ are left exact (and in $\mathcal{M}$ finite $\varprojlim$ are representable – follows from \textup{a} and \textup{c}).
\item[\upshape c)] $\mathcal{M}$ has a cardinal filtration $(\Filt^\pi)_{\pi\geq\pi_0}$.
\item[\upshape d)] The $\Filt^\pi$ are stable under $\varprojlim$ indexed by categories $I$ such that $\operatorname{card}\Fl(I) \leq \pi$, and under quotients ($Z \in \Filt^\pi$ and $Z \longrightarrow Z'$ an epimorphism implies $Z' \in \Filt^\pi$).
\item[\upshape e)] If $A_1$, $A_2$ are two subobjects of an object $B$ of $\mathcal{M}$, and $A_0 = A_1 \cap A_2$, then $A_1 \amalg_{A_0} A_2 \longrightarrow B$ is a monomorphism.
\item[\upshape f)] If we have a cocartesian square
\begin{center}
\begin{tikzcd}
A_0 \arrow[r, "i"] \arrow[d] & A_1 \arrow[d] \\
A_2 \arrow[r, "i'"] & A_3
\end{tikzcd}
\end{center}
with $i$ a monomorphism, then $i'$ is a monomorphism.
\end{itemize}
\end{minipage}\right. \tag*{(2.6.1)} \]
$^{(10,11,12,13)}$\footnotetext[10]{\textbf{N.B.} a + c also means that $\mathcal{M}$ is a ``$\mathfrak{U}$-paratopos category'', and implies the existence of a cardinal filtration satisfying d) (except the assertion concerning the quotients $Z'$ of $Z$). Thus we are dealing with a $\mathfrak{U}$-paratopos satisfying in addition b) e) f). (Is it then a topos\,? \textit{No,} for e), f) are O.K. in \textit{abelian categories\,!})}\footnotetext[11]{\textbf{N.B.} Condition d) is probably unnecessary, at least if epimorphisms are effective, for d) will be satisfied in any case for $\pi = 2^c$, with $c \geq \pi_0$.}\footnotetext[12]{\textbf{N.B.} Condition e) is \textit{not} satisfied for the category (Gr), and the theorem fails for this category, which does \textit{not} have ``enough injectives'', cf. SGA 4 I 8.12.9. On the other hand, e) and f) are satisfied by \textit{abelian} categories.}\footnotetext[13]{Cf. XVIII 4.6.7 (p. 157) for a more congenial formulation of the set of conditions a) to f) for $\mathcal{M}$.}. \textit{(We shall then say that $\mathcal{M}$ is} ``very special''\textit{). Let moreover}
\[ W \ \subset \Fl(\mathcal{M}) \tag*{(2.6.2)} \]
\textit{be a localizer in $\mathcal{M}$, which satisfies the conditions}
\[ \left\{ \begin{minipage}{0.72\textwidth}\itshape
\begin{itemize}[leftmargin=2.5em]
\item[\upshape a)] $W$ is stable under filtered $\varinjlim$ $^{\textup{(14)}}$.
\item[\upshape b)] $W$ is moreover $\mathrm{L}_1$-accessible, i.e. there exists a \textup{small} subset $W_0 \subset W$ such that $W$ is the $\mathrm{L}_1$-envelope of $W$, i.e. such that $W$ is the smallest subset of $\Fl(\mathcal{M})$ containing $W_0$, and stable under filtered $\varinjlim$. % [sic] la L_1-enveloppe de W -> la L_1-enveloppe de W_0, W_0 being the small generating part just introduced
\end{itemize}
\end{minipage}\right. \tag*{(2.6.3)} \]
\footnotetext[14]{\textbf{N.B.} a) already implies that $W = W^\natural$.}

\printedpage{16}
% head: A. Grothendieck : Dérivateurs, Ch. XIX, Édition des Universités de Montpellier II et Paris VII
% foot: Version 11 février 2019

\msp{page 26}
\textit{Let us then put}
\[ C \ = \ \Mono(\mathcal{M}), \qquad TC \ = \ C\cap W\ , \qquad TF \ = \ C_*\ , \qquad F = (TC)_*\ . \tag*{(2.6.4)} \]
\textit{Under these conditions:}
\begin{itemize}[leftmargin=2.5em]
\item[1°)] \textit{$C \Qcouple TF$, more precisely $C \qpair TF$ is a special Quillen couple.} % reading resolved: DER_XIX_corrections.csv
\item[2°)] \textit{For $TC \qpair F$ to be a Quillen couple, it is necessary and sufficient that $TC$ be stable under cobase change, and then it is a special Quillen couple.} % reading resolved: DER_XIX_corrections.csv
\item[3°)] \textit{For $(W,C,F)$ to be a Quillen triple (which will necessarily be special by \textup{1°, 2°}, and in particular closed), it is necessary and sufficient that the following conditions be satisfied.}
\[ \left\{ \begin{array}{ll}
\text{a)} & TC\ (\deq W\cap\Mono)\ \text{\itshape is stable under cobase change.}\\[1ex]
\text{b)} & TF\ (\deq (\Mono)_*) \ = \ \text{\itshape ``maps of injective type''}\ \subset W.
\end{array}\right. \tag*{(2.6.5)~$^{(15)}$} \]
\footnotetext[15]{\textbf{N.B.} These conditions (2.6.5) are equivalent to saying that $(W,C)$ is a structure of a category with cofibrations on $\mathcal{M}$.}
\textit{And condition} b) \textit{is also equivalent to:}
\[ \begin{minipage}{0.7\textwidth}
\begin{itemize}[leftmargin=2.5em]
\item[b$'$)] \textit{Every arrow $f$ of $\mathcal{M}$ (or only of $TF$) factors as $f = pi$, with $i$ a monomorphism, $p \in W$.}
\end{itemize}
\end{minipage} \tag*{(2.6.6)} \]
\end{itemize}

\textbf{Corollary 2.7.} \textit{$\mathcal{M}$ and $W$ being as in} (2.6.1), (2.6.3), \textit{for $(W,C,F)$ to be a Quillen triple that is} left proper \textit{(i.e. $C \subset \Cof_W$), it is necessary and sufficient that $W$ satisfy condition} (2.6.5. b) \textit{(i.e. $W$ contains the arrows of injective type), and moreover}
\msp{page 27}

\begin{itemize}[leftmargin=7em]
\item[] $C \subset \Cof_W$, \textit{i.e. the condition:}
\item[a$'$)] \textit{For every cocartesian diagram}
\[ \begin{tikzcd}
A_0 \arrow[r, "i"] \arrow[d, "f"'] & A_1 \arrow[d, "f'"'] \\
A_2 \arrow[r, "i'"] & A_3
\end{tikzcd} \tag*{(2.7.7)} \]
\textit{with $i$ a monomorphism, if $f \in W$, then $f' \in W$. (Stability of $W$ under cobase change along monomorphisms.)}
\end{itemize}

\textit{Under these conditions, this Quillen triple is special, and a fortiori closed.}

This corollary follows from 2.6. 3°), since $C \subset \Cof_W$ implies $TC \deq C\cap W = C\cap W^{\mathrm{univ}}$ (for $\Cof_W \cap W \ = \ W^{\mathrm{univ}}$), hence that $TC$ is stable under cobase change, whence condition (2.6.5. a).

\printedpage{17}
% head: A. Grothendieck : Dérivateurs, Ch. XIX, Édition des Universités de Montpellier II et Paris VII
% foot: Version 11 février 2019

\textbf{2.8.} We say that a \textit{Quillen structure} $(W,C,F)$ on $\mathcal{M}$ is \textit{very special,} if it is obtained as in theorem 2.1., i.e. if we have the conditions:
\[ \left\{ \begin{minipage}{0.72\textwidth}
\begin{itemize}[leftmargin=2.5em]
\item[a)] $\mathcal{M}$ is very special, i.e. satisfies conditions (2.6.1).
\item[b)] $W$ is stable under small filtered $\varinjlim$, and is accessible (conditions (2.6.3)).
\item[c)] $C = \Mono$.
\end{itemize}
\end{minipage}\right. \tag*{(2.8.1)} \]
Thus for a localizer $W$ in $\mathcal{M}$ to correspond to a very special Quillen structure \msp{page 28} in $\mathcal{M}$ (which will necessarily be unique), it is necessary and sufficient that $\mathcal{M}$ and $W$ satisfy conditions (2.8.1. a)\,b)), and that moreover we have:
\[ \left\{ \begin{minipage}{0.65\textwidth}
\begin{itemize}[leftmargin=2.5em]
\item[a)] $TC \deq \Mono\cap W$ is stable under cobase change.
\item[b)] $TF \deq (\Mono)_* \subset W$, i.e. $W$ contains the arrows of injective type in $\mathcal{M}$.
\end{itemize}
\end{minipage}\right. \tag*{(2.8.2)} \]
\footnote[16]{We also say that the localizer $W$ is very special, if it defines a very special Quillen structure.}.

We now find the

\textbf{Theorem 2.9.} \textit{Let} $(W, C = \Mono, F = (TC)_* = (W\cap\Mono)_*)$ \textit{be a very special Quillen triple in a category $\mathcal{M}$} (2.8). \textit{Then:}
\begin{itemize}[leftmargin=2.5em]
\item[1°)] \textit{For every $I$ in} Cat, $(W(I),C(I),F[I])$ \textit{is a very special Quillen triple in $\mathcal{M}(I)$, and in particular we have}
\[ C(I) \ = \ \Mono(\mathcal{M}(I))\ , \tag*{(2.9.1)} \]
\textit{and moreover}
\[ F[I]\cap W(I) \ = \ (TF)[I]\ . \tag*{(2.9.2)} \]
\item[2°)] \textit{For every arrow $u : I \longrightarrow I'$ in} Cat, \textit{whence the pair of adjoint functors}
\[ \begin{tikzcd}
\mathcal{M}(I) \arrow[r, shift right, "u_*"'] & \mathcal{M}(I') \arrow[l, shift right, "u^*"']
\end{tikzcd}\ , \tag*{(2.9.3)} \]
\textit{we have}
\[ \left\{ \begin{array}{rcl}
u^*(W(I')) & \subset & W(I)\\
u^*(C(I')) & \subset & C(I)\\
u^*(TC(I')) & \subset & TC(I)\\
u_*(F[I]) & \subset & F[I']\\
u_*(TF[I]) & \subset & TF[I']\ ,
\end{array}\right. \tag*{(2.9.4)} \]
\textit{finally $u_*$ induces a functor}
\printedpage{18}
% head: A. Grothendieck : Dérivateurs, Ch. XIX, Édition des Universités de Montpellier II et Paris VII
% foot: Version 11 février 2019
\msp{page 29}
\[ u_*^{\mathrm{f}} \ :\ \mathcal{M}(I)_{\mathrm{f}} \longrightarrow \mathcal{M}(I')_{\mathrm{f}} \qquad\quad ^{(17)} \tag*{(2.9.5)} \]
\footnotetext[17]{full subcategories of fibrant objects.}
compatible with the localizers $W(I)_{\mathrm{f}}$ \textit{and} $W(I')_{\mathrm{f}}$ \textit{induced by $W(I)$, $W(I')$ respectively, hence induces a functor}
\[ u_*^W \ :\ \Ho_W(I) \longrightarrow \mathit{Ho}_W(I')\ . \tag*{(2.9.6)} \] % [sic] italic Ho_W(I') -> upright Ho_W(I'), inconsistent with the upright Ho_W elsewhere
\textit{This functor is a} right adjoint \textit{to the functor}
\[ u_W^* \ :\ \mathit{Ho}_W(I') \longrightarrow \Ho_W(I) \tag*{(2.9.7)} \] % [sic] italic Ho_W(I') -> upright Ho_W(I'), inconsistent with the upright Ho_W elsewhere
\textit{induced by $u^*$} (2.9.3).
\end{itemize}
\textsc{Proof of 2.9.} To prove 1°), we must first check that the preliminary conditions (2.8.1. a)\,b)) on $\mathcal{M}$, $W$ are stable under passage from $\mathcal{M}$, $W$ to $\mathcal{M}(I)$, $W(I)$, which is done by AQT \texttt{[âne qui trotte: at a donkey's trot, i.e. by routine verification]}. Besides, one checks right away that
\[ C(I) \ = \ \Mono(\mathcal{M}(I))\ , \]
and theorem 2.6 (1°, 2°) then implies that
\[ \begin{array}{ccccc}
C(I) & \qpair & C(I)_* & \deq & TF[I]\\
\underbrace{TC(I)}_{\deq C(I)\cap W(I)} & \qpair & (TC(I))_* & \deq & F[I]
\end{array} \] % reading resolved: DER_XIX_corrections.csv
are special Quillen couples. (Noting for the second, that $TC(I)$ is indeed \msp{page 30} stable under cobase change, $TC$ being so.) To check that we have a very special Quillen triple, it remains by virtue of 2.8 to check that
\[ TF[I] \ \subset\ W(I)\ , \]
which is written
\[ v^*(TF[I]) \ \subset\ W(I_0)\ , \]
where
\[ v \ :\ I_0 \longrightarrow I \]
is defined as in (1.12) (p. 10). In fact, we shall even have
\[ v^*(TF[I]) \ = \ v^*v_*(TF(I_0)) \ \subset\ \underbrace{TF(I_0)}_{= F(I_0)\cap W(I_0)}\quad , \]
and the proof is dual to that of
\[ v^*v_!(TC(I_0)) \ \subset\ TC(I_0) \]
(cf. (2.3.7), p. 21), it does not depend on questions of factorization in $\mathcal{M}(I)$, and comes here from the stability of $TF$ under small products. This proves 1°.

The proof of 2° is now the dual of that of (2.3. 2°). Since this makes us repeat the same proof twice, this shows that we have not ``set up the sorites'' quite carefully enough. So by way of making up for this, I add a